\chapter{HIỆN THỰC HÓA HỆ THỐNG}

\section{Tổng quan cấu trúc mã nguồn thực tế}

Hệ thống MarketFlow được tổ chức thành ba khối mã nguồn độc lập gồm \texttt{backend/} (API nghiệp vụ và máy chủ MCP), \texttt{ai-runner/} (dịch vụ điều phối phiên làm việc AI cô lập), và \texttt{frontend/} (ứng dụng web một trang SPA). Bảng~\ref{tab:source-tree} trình bày chi tiết cấu trúc phân bổ các thành phần trong dự án.

\begin{table}[H]
\centering
\footnotesize
\renewcommand{\arraystretch}{1.05}
\caption{Cấu trúc phân bổ mã nguồn hệ thống MarketFlow}
\label{tab:source-tree}
\begin{tabularx}{\textwidth}{|L{4.2cm}|Y|L{3.1cm}|}
\hline
\thead{Đường dẫn thư mục} & \thead{Mô tả thành phần mã nguồn} & \thead{Công nghệ sử dụng}\\
\hline
\texttt{backend/app/core} & Cấu hình môi trường (\texttt{config.py}), xử lý ngoại lệ chuẩn hóa và bảo mật (\texttt{security.py}: Bcrypt, JWT, phân quyền vai trò). & Python 3.12, Passlib, PyJWT\\
\hline
\texttt{backend/app/db} & Quản lý phiên kết nối cơ sở dữ liệu (\texttt{session.py}), cấu trúc ORM cơ sở và kịch bản khởi tạo dữ liệu mẫu (\texttt{seed.py}). & SQLAlchemy 2.0, PostgreSQL / SQLite\\
\hline
\texttt{backend/app/models} & Định nghĩa các bảng dữ liệu quan hệ (User, Workspace, Campaign, ContentItem, Schedule, MetricEntry, AiThread, AiAction, AuditLog). & SQLAlchemy 2.0 Declarative\\
\hline
\texttt{backend/app/schemas} & Đặc tả hợp đồng dữ liệu đầu vào/đầu ra và quy tắc kiểm tra tính hợp lệ. & Pydantic v2\\
\hline
\texttt{backend/app/api/v1} & Các bộ định tuyến REST API: xác thực, chiến dịch, nội dung, phê duyệt, lịch đăng, số liệu KPI và nhật ký. & FastAPI Routers\\
\hline
\texttt{backend/app/assistant} & Máy chủ công cụ ngữ cảnh (\texttt{mcp\_server.py}), quản lý luồng hội thoại (\texttt{chat\_service.py}) và hàng đợi đề xuất hành động (\texttt{actions\_service.py}). & FastMCP, SSE Streaming\\
\hline
\texttt{ai-runner/app} & Dịch vụ điều phối tiến trình AI (\texttt{runner.py}, \texttt{events.py}): quản lý hàng đợi, cắt chuỗi thời gian thực và xoay vòng đa mô hình dự phòng. & FastAPI, Asyncio, OpenCode CLI\\
\hline
\texttt{frontend/src/features} & Các phân trang nghiệp vụ theo miền: \texttt{dashboard}, \texttt{campaigns}, \texttt{contents}, \texttt{schedules}, \texttt{metrics}, \texttt{assistant}, \texttt{brand}, \texttt{audit}. & React 18, TypeScript, TanStack Query\\
\hline
\texttt{frontend/src/components} & Hệ thống khung bố cục 3 cột (\texttt{AppShell}), ngăn kéo trợ lý AI (\texttt{AssistantDrawer}), thẻ đề xuất hành động (\texttt{ActionProposalCard}) và bảng dữ liệu. & Tailwind CSS, Lucide Icons\\
\hline
\end{tabularx}
\end{table}

\section{Hiện thực hóa các phân hệ Backend}

\subsection{Phân hệ Bảo mật, Xác thực và Phân quyền mức bản ghi (Record-level Scope)}

Phân hệ an ninh trong \texttt{backend/app/core/security.py} kết hợp kiểm tra đồng thời ba điều kiện trên mỗi yêu cầu truy cập: chữ ký số JWT hợp lệ, trạng thái tài khoản đang hoạt động trong cơ sở dữ liệu (\texttt{user.is\_active == True}), và quyền hạn của người dùng trong không gian làm việc hiện tại.

Nhằm ngăn chặn lỗi truy cập trái phép tài nguyên ngang hàng (IDOR -- Insecure Direct Object Reference), hệ thống áp dụng cơ chế kiểm soát phạm vi bản ghi ngay tại tầng truy vấn dữ liệu. Quyền thao tác của từng nhóm người dùng được quy định cụ thể trong Bảng~\ref{tab:scope-matrix}.

\begin{table}[H]
\centering
\small
\caption{Ma trận phân quyền truy cập tài nguyên theo vai trò nghiệp vụ}
\label{tab:scope-matrix}
\begin{tabularx}{\textwidth}{|L{2.2cm}|L{2.8cm}|C{2.6cm}|Y|C{2.0cm}|}
\hline
\thead{Tài nguyên} & \thead{Hành động} & \thead{Quản lý (Manager)} & \thead{Nhân viên sở hữu / Được phân công} & \thead{Nhân viên ngoài phạm vi}\\
\hline
\multirow{3}{2.2cm}{\raggedright Chiến dịch} & Xem / Tìm kiếm & Toàn quyền & Xem chiến dịch trong không gian làm việc & Từ chối (HTTP 403)\\
\cline{2-5}
& Tạo mới / Sửa & Toàn quyền & Sửa chiến dịch được phân công & Từ chối (HTTP 403)\\
\cline{2-5}
& Xóa chiến dịch & Toàn quyền & Không có quyền xóa & Từ chối (HTTP 403)\\
\hline
\multirow{2}{2.2cm}{\raggedright Nội dung} & Soạn / Sửa nháp & Toàn quyền & Soạn trong chiến dịch được giao & Từ chối (HTTP 403)\\
\cline{2-5}
& Phê duyệt (Approve) & Toàn quyền & Không có quyền duyệt & Từ chối (HTTP 403)\\
\hline
Chỉ số kênh & Nhập / Xuất báo cáo & Toàn quyền & Nhập cho chiến dịch được giao & Từ chối (HTTP 403)\\
\hline
\end{tabularx}
\end{table}

\subsection{Hiện thực hóa Máy trạng thái duyệt nội dung (HITL State Machine)}

Để đảm bảo quy trình kiểm duyệt của con người (Human-in-the-loop), phân hệ quản lý nội dung áp dụng các điều kiện kiểm soát chuyển trạng thái chặt chẽ tại tầng dịch vụ. Bảng~\ref{tab:hitl-rules} tổng hợp các quy tắc chuyển trạng thái bài viết và cơ chế bảo vệ dữ liệu sau khi duyệt.

\begin{table}[H]
\centering
\small
\caption{Quy tắc chuyển trạng thái nội dung và kiểm soát phê duyệt}
\label{tab:hitl-rules}
\begin{tabularx}{\textwidth}{|L{2.2cm}|L{2.6cm}|C{2.2cm}|Y|C{2.2cm}|}
\hline
\thead{Trạng thái gốc} & \thead{Thao tác kích hoạt} & \thead{Quyền hạn} & \thead{Điều kiện bảo vệ (Guard Condition)} & \thead{Trạng thái đích}\\
\hline
Khởi tạo mới & \texttt{POST /contents} & Marketer, Manager & Bài viết mới hoặc đề xuất từ AI luôn khởi tạo ở trạng thái nháp; không cho phép tạo trực tiếp bài đã duyệt. & \texttt{DRAFT} hoặc \texttt{AI\_DRAFT}\\
\hline
\texttt{DRAFT} & \texttt{POST .../submit} & Tác giả bài viết & Tiêu đề và nội dung bài viết không được để trống; thuộc chiến dịch đang hoạt động. & \texttt{IN\_REVIEW}\\
\hline
\texttt{IN\_REVIEW} & \texttt{POST .../approve} & Manager & Chỉ người dùng có vai trò Quản lý mới được phép phê duyệt và ghi nhận vào nhật ký kiểm duyệt. & \texttt{APPROVED}\\
\hline
\texttt{IN\_REVIEW} & \texttt{POST .../reject} & Manager & Bắt buộc đính kèm nhận xét hoặc lý do từ chối để nhân viên chỉnh sửa lại. & \texttt{REJECTED}\\
\hline
\texttt{APPROVED} & \texttt{PUT /contents/\{id\}} & Tác giả, Manager & \textbf{Tự động thu hồi trạng thái duyệt}: Khi nội dung đã duyệt bị chỉnh sửa, hệ thống chuyển bài viết về trạng thái nháp để thẩm định lại. & \texttt{DRAFT}\\
\hline
\texttt{APPROVED} & \texttt{POST /schedules} & Marketer, Manager & Chỉ các bài viết ở trạng thái \texttt{APPROVED} mới được phép lên lịch xuất bản. & \texttt{SCHEDULED}\\
\hline
\end{tabularx}
\end{table}

\section{Hiện thực hóa giao diện làm việc 3 cột (Split-Pane Workspace)}

Giao diện người dùng được thiết kế theo hệ màu sáng dịu mắt (\textit{Soft Daylight}: nền tổng thể \texttt{\#F6F7F9}, bề mặt bảng \texttt{\#FFFFFF}, vùng làm việc AI \texttt{\#F0F4F8}) kết hợp bộ phông chữ chuẩn tiếng Việt (\texttt{Be Vietnam Pro}, \texttt{Plus Jakarta Sans} và số liệu đẳng rộng \texttt{JetBrains Mono} \texttt{tabular-nums}) giúp người dùng làm việc liên tục nhiều giờ với bảng số liệu mà không gây mỏi thị giác.

\begin{figure}[H]
\centering
\resizebox{0.98\textwidth}{!}{%
\begin{tikzpicture}[
 font=\sffamily\scriptsize,
 >=Latex,
 frame/.style={draw=black!80,fill=white,rounded corners=3pt,line width=0.8pt},
 panel/.style={draw=black!65,fill=white,rounded corners=2pt},
 card/.style={draw=black!60,fill=black!3,rounded corners=2pt},
 btn/.style={draw=black!85,fill=black!6,rounded corners=2pt,inner sep=3pt,font=\sffamily\bfseries\scriptsize},
 btn_dark/.style={draw=black,fill=black!16,text=black,rounded corners=2pt,inner sep=3pt,font=\sffamily\bfseries\scriptsize}
]

\draw[draw=black!85,fill=white,rounded corners=4pt,line width=0.9pt] (0, 0) rectangle (16.0, 8.6);

\draw[draw=black!85,fill=black!10,rounded corners=3pt] (0, 7.80) rectangle (16.0, 8.6);
\node[font=\sffamily\bfseries\scriptsize,text=black,anchor=west] at (0.35, 8.20) {KHÔNG GIAN LÀM VIỆC 3 CỘT (SPLIT-PANE WORKSPACE) -- MARKETFLOW};
\node[font=\sffamily\scriptsize,text=black,anchor=east] at (15.65, 8.20) {Tài khoản: Vũ Hiếu Kiên (Quản lý)};

\draw[frame] (0.25, 0.25) rectangle (3.45, 7.55);
\node[font=\sffamily\bfseries\scriptsize,text=black,anchor=west] at (0.45, 7.15) {ĐIỀU HƯỚNG CHÍNH};
\node[anchor=west,font=\sffamily\scriptsize] at (0.45, 6.60) {1. Tổng quan};
\node[anchor=west,font=\sffamily\bfseries\scriptsize] at (0.45, 6.08) {2. Chiến dịch};
\node[anchor=west,font=\sffamily\scriptsize] at (0.45, 5.56) {3. Nội dung \& Duyệt bài};
\node[anchor=west,font=\sffamily\scriptsize] at (0.45, 5.04) {4. Lịch xuất bản};
\node[anchor=west,font=\sffamily\scriptsize] at (0.45, 4.52) {5. Số liệu \& Báo cáo KPI};
\node[anchor=west,font=\sffamily\scriptsize] at (0.45, 4.00) {6. Trợ lý Marketing AI};
\node[anchor=west,font=\sffamily\scriptsize] at (0.45, 3.48) {7. Tài sản thương hiệu};
\node[anchor=west,font=\sffamily\scriptsize] at (0.45, 2.96) {8. Nhật ký hệ thống};
\draw[draw=black!35,line width=0.5pt] (0.45, 2.50) -- (3.25, 2.50);
\node[anchor=west,font=\sffamily\scriptsize] at (0.45, 2.05) {Giao diện: Soft Daylight};
\node[anchor=west,font=\sffamily\scriptsize] at (0.45, 1.55) {Backend: FastAPI};
\node[anchor=west,font=\sffamily\scriptsize] at (0.45, 1.05) {AI: Đa mô hình tự động};

\node[font=\sffamily\bfseries\scriptsize,text=black,anchor=west] at (3.70, 7.25) {Chiến dịch: Khuyến mãi Mùa Hè 2026};
\node[btn,anchor=east] at (12.40, 7.25) {Xuất báo cáo Excel};
\node[btn_dark,anchor=east] at (15.75, 7.25) {+ Tạo chiến dịch mới};

\node[card, minimum width=2.75cm, minimum height=1.10cm, align=left, inner sep=4pt] at (5.10, 6.25) {
 {\scriptsize Tổng ngân sách}\\[2pt]
 \textbf{\scriptsize 48M / 60M VNĐ}
};

\node[card, minimum width=2.75cm, minimum height=1.10cm, align=left, inner sep=4pt] at (8.10, 6.25) {
 {\scriptsize Chi phí / Chuyển đổi}\\[2pt]
 \textbf{\scriptsize 95.000đ (CPA)}
};

\node[card, minimum width=2.75cm, minimum height=1.10cm, align=left, inner sep=4pt] at (11.10, 6.25) {
 {\scriptsize Lượt hiển thị}\\[2pt]
 \textbf{\scriptsize 124.500 lượt}
};

\node[card, minimum width=2.75cm, minimum height=1.10cm, align=left, inner sep=4pt] at (14.20, 6.25) {
 {\scriptsize Tỷ lệ nhấp (CTR)}\\[2pt]
 \textbf{\scriptsize 3.42\% (+0.6\%)}
};

\draw[panel] (3.70, 0.25) rectangle (9.65, 5.45);
\node[font=\sffamily\bfseries\scriptsize,text=black,anchor=west] at (3.90, 5.10) {BẢNG DỮ LIỆU CHIẾN DỊCH \& NỘI DUNG};
\draw[draw=black!35,line width=0.5pt] (3.70, 4.82) -- (9.65, 4.82);

\node[font=\sffamily\scriptsize,anchor=west,text width=5.5cm] at (3.90, 4.42) {\textbf{1. Ưu đãi giải nhiệt mùa hè}};
\node[font=\sffamily\scriptsize,anchor=west,text width=5.5cm] at (3.90, 4.02) {Kênh: Facebook -- Trạng thái: Chờ duyệt};
\draw[draw=black!20,line width=0.4pt] (3.90, 3.70) -- (9.45, 3.70);

\node[font=\sffamily\scriptsize,anchor=west,text width=5.5cm] at (3.90, 3.32) {\textbf{2. Kịch bản video ngắn máy lọc ProAir}};
\node[font=\sffamily\scriptsize,anchor=west,text width=5.5cm] at (3.90, 2.92) {Kênh: TikTok -- Trạng thái: Bản nháp AI};
\draw[draw=black!20,line width=0.4pt] (3.90, 2.60) -- (9.45, 2.60);

\node[font=\sffamily\scriptsize,anchor=west,text width=5.5cm] at (3.90, 2.22) {\textbf{3. Thư điện tử chăm sóc khách hàng}};
\node[font=\sffamily\scriptsize,anchor=west,text width=5.5cm] at (3.90, 1.82) {Kênh: Email -- Trạng thái: Đã phê duyệt};
\draw[draw=black!20,line width=0.4pt] (3.90, 1.50) -- (9.45, 1.50);

\node[font=\sffamily\scriptsize,anchor=west,text width=5.5cm] at (3.90, 1.10) {\textbf{4. Bài giới thiệu màng lọc HEPA 3 lớp}};
\node[font=\sffamily\scriptsize,anchor=west,text width=5.5cm] at (3.90, 0.70) {Kênh: Zalo OA -- Trạng thái: Đã lên lịch};

\draw[panel,fill=black!2] (9.90, 0.25) rectangle (15.75, 5.45);
\node[font=\sffamily\bfseries\scriptsize,text=black,anchor=west] at (10.10, 5.10) {KHUNG TÁC VỤ AI (CONTEXTUAL AI PANEL)};
\draw[draw=black!35,line width=0.5pt] (9.90, 4.82) -- (15.75, 4.82);

\node[font=\sffamily\scriptsize,anchor=west,text width=5.4cm] at (10.10, 4.42) {Chiến dịch: Mùa Hè 2026 | Kênh: Facebook};
\node[font=\sffamily\scriptsize,anchor=west,text width=5.4cm] at (10.10, 4.02) {Kiểm tra Brand Kit: Hợp lệ (0 từ cấm)};

\draw[draw=black!50,fill=white,rounded corners=2pt] (10.10, 1.32) rectangle (15.55, 3.72);
\node[font=\sffamily\scriptsize,anchor=north west,text width=5.1cm] at (10.25, 3.58) {\textbf{Đề xuất bản nháp AI (\#42):}};
\node[font=\sffamily\scriptsize,anchor=north west,text width=5.1cm] at (10.25, 3.10) {
 \textit{``Không khí sạch cho cả gia đình! Khám phá máy lọc ProAir màng HEPA 3 lớp.''}
};
\node[font=\sffamily\scriptsize,anchor=south west,text width=5.1cm] at (10.25, 1.45) {Trạng thái: PROPOSED | Chờ xác nhận};

\node[btn_dark,anchor=west] at (10.10, 0.75) {Áp dụng vào bảng};
\node[btn,anchor=west] at (12.60, 0.75) {Chỉnh sửa};
\node[btn,anchor=east] at (15.55, 0.75) {Huỷ bỏ};

\end{tikzpicture}%
}

\caption{Kiến trúc không gian làm việc 3 cột (Split-Pane Workspace) và khung tác vụ AI của MarketFlow}
\label{fig:ui-wireframe}
\end{figure}

\subsection{Bố cục không gian làm việc và các phân trang chính}
\begin{enumerate}
 \item \textbf{Cột điều hướng thu gọn bên trái (Collapsible Sidebar)}: Cho phép chuyển đổi nhanh giữa 8 phân trang nghiệp vụ (Tổng quan, Chiến dịch, Nội dung, Lịch đăng, Số liệu \& KPI, Trợ lý AI, Tài sản thương hiệu, Nhật ký hệ thống) mà không cần tải lại trang.
 \item \textbf{Khu vực làm việc trung tâm (Main Workspace)}: Hiển thị hàng thẻ chỉ số KPI trọng yếu ở phía trên, biểu đồ phân tích hiệu quả đa kênh và bảng dữ liệu có cố định tiêu đề cột khi cuộn. Tích hợp sẵn chức năng \textbf{Xuất báo cáo Excel (CSV UTF-8 BOM)} và nút gửi ngữ cảnh nhanh sang AI trên từng dòng chiến dịch.
 \item \textbf{Khung tác vụ AI ngữ cảnh bên phải (Contextual AI Copilot Drawer)}: Cho phép vừa theo dõi bảng số liệu ở giữa vừa trao đổi với Trợ lý AI ở cột bên phải. Luồng phản hồi SSE được duy trì ở cấp module giúp người dùng chuyển trang hoặc đổi tab trình duyệt mà không bị ngắt phiên hội thoại đang chạy.
 \item \textbf{Thẻ xác nhận đề xuất hành động (Action Proposal Card)}: Khi AI đề xuất tạo bản nháp mới, lập lịch đăng hoặc điều chỉnh chiến dịch qua các công cụ \texttt{propose\_*}, hệ thống hiển thị thẻ tóm tắt trực quan để người dùng xem trước và bấm \textit{Áp dụng vào bảng} hoặc \textit{Huỷ bỏ}.
\end{enumerate}

\section{Hiện thực hóa Cơ chế Điều phối Đa mô hình và Chống ảo giác số liệu}

Để đảm bảo tốc độ phản hồi nhanh cùng tính sẵn sàng liên tục, dịch vụ \texttt{ai-runner} (\texttt{ai-runner/app/runner.py}) kết hợp với máy chủ công cụ MCP (\texttt{backend/app/assistant/mcp\_server.py}) triển khai cơ chế xoay vòng đa nhà cung cấp và kiểm soát trung thực dữ liệu:
\begin{align}
\text{Total\_Views} &= \sum_{m \in M} m.\text{views}, \quad \text{Total\_Clicks} = \sum_{m \in M} m.\text{clicks} \\
\text{Total\_Cost} &= \sum_{m \in M} m.\text{cost}, \quad \text{Total\_Conversions} = \sum_{m \in M} m.\text{conversions} \\
\text{CTR} &= \begin{cases} \dfrac{\text{Total\_Clicks}}{\text{Total\_Views}} \times 100\% & \text{khi } \text{Total\_Views} > 0 \\ 0.0\% & \text{khi } \text{Total\_Views} = 0 \end{cases} \\
\text{CPA} &= \begin{cases} \dfrac{\text{Total\_Cost}}{\text{Total\_Conversions}} & \text{khi } \text{Total\_Conversions} > 0 \\ 0 & \text{khi } \text{Total\_Conversions} = 0 \end{cases}
\end{align}

Hệ thống áp dụng ba nguyên tắc kiểm soát chất lượng phản hồi AI:
\begin{enumerate}
 \item \textbf{Tính toán chỉ số xác định từ cơ sở dữ liệu}: Các chỉ số CTR, CVR, CPC, CPA và tiến độ giải ngân được tính toán chính xác bằng mã nguồn Python trong các công cụ \texttt{marketflow\_*} trước khi cung cấp cho mô hình ngôn ngữ; mô hình không tự ước lượng số liệu.
 \item \textbf{Tự động chuyển đổi dự phòng đa nhà cung cấp (Multi-Provider Fallback)}: Hệ thống kết nối đồng thời OpenCode Zen, Groq LPU, Google Gemini và OpenRouter. Khi một mô hình chạm giới hạn tần suất (HTTP 429) hoặc quá thời gian chờ, bộ điều phối tự động đánh dấu thời gian nghỉ (\textit{per-model cooldown}) và chuyển yêu cầu sang mô hình dự phòng kế tiếp.
 \item \textbf{Cô lập quyền ghi qua cơ chế Đề xuất (Proposal-Only)}: Trợ lý AI chỉ có quyền đọc dữ liệu trong phạm vi không gian làm việc và tạo bản ghi đề xuất ở trạng thái \texttt{PROPOSED}; mọi thao tác ghi chính thức vào bảng nghiệp vụ đều yêu cầu xác nhận từ người dùng.
\end{enumerate}

\section{Kết luận chương}

Chương 5 đã trình bày chi tiết quá trình hiện thực hóa hệ thống MarketFlow: từ cơ chế phân quyền theo phạm vi bản ghi, máy trạng thái phê duyệt nội dung, thiết kế không gian làm việc 3 cột dịu mắt đến kiến trúc điều phối trợ lý AI đa mô hình có kiểm soát đề xuất. Các thành phần này là cơ sở thực tế để tiến hành kiểm thử tự động và đánh giá hiệu năng ở Chương 6 và Chương 7.
